\section*{Chapitre \ref{ch:qm:monte-carlo} --- Monte-Carlo}

\begin{solution}{exo:qm:monte-carlo:1}
L'erreur varie comme $n^{-1/2}$~: un facteur quatre exige $16 \times 100\,000 = 1.6$ million de trajectoires, soit un facteur de \omterm{def:qm:monte-carlo:mc}{réduction de variance} de 16 à 100\,000 trajectoires.
\end{solution}

\begin{solution}{exo:qm:monte-carlo:2}
$1/(1 - \rho^2)$~: 5.3 pour $\rho = 0.9$ et 50.3 pour $\rho = 0.99$. De fortes corrélations sont nécessaires pour obtenir de grands facteurs.
\end{solution}

\begin{solution}{exo:qm:monte-carlo:3}
$\Var(\frac12(f(Z) + f(-Z))) = \frac12\Var f(Z)(1 + \mathrm{Corr}(f(Z), f(-Z)))$ par paire, c'est-à-dire par deux évaluations~; cet \omterm{def:qm:estimation:estimator}{estimateur} bat l'échantillonnage simple lorsque la corrélation est négative, ce qu'assure une fonction $f$
monotone. Pour une fonction paire $f(Z) = f(-Z)$, la corrélation vaut 1 et la paire équivaut à un seul tirage pour le coût de deux.
\end{solution}

\begin{solution}{exo:qm:monte-carlo:4}
Par la formule de la variance totale, $\sigma^2 = \sum_sp_s\sigma_s^2 + \sum_sp_s(\mu_s - \mu)^2$. L'\omterm{def:qm:estimation:estimator}{estimateur} stratifié $\sum_sp_s\bar Y_s$ avec $n_s = np_s$ a pour variance $\sum_sp_s^2\sigma_s^2/n_s =
\frac1n\sum_sp_s\sigma_s^2$, ce qui supprime le terme inter-strates, par nature positif ou nul.
\end{solution}

\begin{solution}{exo:qm:monte-carlo:5}
L'\omterm{def:qm:estimation:estimator}{estimateur} est $\mathbf 1\{Z > z^*\}e^{-\theta Z + \theta^2/2}$ avec $Z \sim \mathcal N(\theta, 1)$~; son moment d'ordre deux vaut
\[\E_f\bigl[\mathbf 1\{Z > z^*\}e^{-\theta Z + \theta^2/2}\bigr] = e^{\theta^2}\bar\Phi(z^* + \theta).\]
La dérivée de son logarithme vaut $2\theta - \varphi(z^* + \theta)/\bar\Phi(z^* + \theta) \approx 2\theta - (z^* + \theta)$ pour de grands arguments (rapport de Mills), qui s'annule en $\theta \approx z^*$. Pour $z^* = 3.67$
le minimiseur exact est 3.80, là où se situe le minimum de la simulation.
\end{solution}

\begin{solution}{exo:qm:monte-carlo:6}
Avec $a(x) = \mu x$ et $b(x) = \sigma x$, $bb' = \sigma^2x$~: $X_{k+1} = X_k(1 + \mu h + \sigma\Delta W + \frac12\sigma^2(\Delta W^2 - h))$. En développant le pas exact,
$e^{(\mu - \sigma^2/2)h + \sigma\Delta W} = 1 + \sigma\Delta W + (\mu - \frac12\sigma^2)h + \frac12\sigma^2\Delta W^2 + O(h^{3/2})$, les mêmes termes.
\end{solution}

\begin{solution}{exo:qm:monte-carlo:7}
Avec $2^{14}$ trajectoires~: corrélation 0.99953, 0.99941 et 0.99425, facteurs de \omterm{def:qm:monte-carlo:mc}{réduction de variance} 1\,071, 853 et 87 pour des prix d'exercice de 80, 100 et 130. Très loin hors de la monnaie, la moyenne géométrique (toujours inférieure à la
moyenne arithmétique) finit dans la monnaie moins souvent que cette dernière, et les deux payoffs se découplent au voisinage du prix d'exercice~: la \omterm{def:qm:monte-carlo:vr}{variable de contrôle} vaut moins là où les deux sont rarement non nuls.
\end{solution}

\begin{solution}{exo:qm:monte-carlo:8}
Les points de Sobol ne sont pas indépendants, donc $s/\sqrt n$ n'est pas l'\omterm{def:qm:estimation:estimator}{erreur type} de leur moyenne~: il mesure la dispersion des payoffs et non celle de l'estimation, et ne dit rien de l'erreur, qui
peut être plus petite ou, avec une mauvaise construction des trajectoires, plus grande. L'erreur doit provenir de brouillages indépendants (répliques), et les points non brouillés ne donnent aucune estimation de l'erreur.
\end{solution}

\begin{solution}{pb:qm:monte-carlo:1}
\textbf{1.} 6.495 avec une \omterm{def:qm:estimation:estimator}{erreur type} de 0.077 pour $2^{14}$ trajectoires~; référence 6.4788 (\omterm{def:qm:monte-carlo:vr}{variable de contrôle} plus Sobol brouillé, 32 répliques de $2^{14}$, \omterm{def:qm:estimation:estimator}{erreur type} 0.00007).
\textbf{2.} 16.
\textbf{3.} À précision égale, le coût vaut $c\sigma^2/\epsilon^2$, donc les \omterm{def:qm:estimation:estimator}{estimateurs} se comparent par le produit de la variance et du coût par tirage~; tous les facteurs donnés ici sont calculés à nombre égal d'évaluations du payoff.
\textbf{4.} Les nombres de chaque trajectoire sont une fonction de (graine, transaction, trajectoire, bloc), et non d'un état séquentiel, de sorte que toute répartition des trajectoires les reproduit.
\textbf{5.} La \omterm{def:qm:monte-carlo:prng}{méthode d'inversion} associe de façon monotone une coordonnée uniforme à une coordonnée normale, de sorte que la structure de l'ensemble de points est préservée~; Box--Muller mélange deux coordonnées au moyen d'un cosinus.
\textbf{6.} 1.8~: les deux payoffs d'une paire ont une corrélation de $-0.43$~; le payoff est monotone dans la trajectoire mais convexe, et sa partie positive élimine les trajectoires baissières.
\textbf{7.} Le call sur moyenne géométrique, 6.1793 en forme fermée~; corrélation 0.99941~; facteur 853.
\textbf{8.} $\beta^* = \Cov(Y, C)/\Var C$~; estimé à partir des mêmes trajectoires que 1.04 (le biais ainsi introduit est d'ordre $1/n$).
\textbf{9.} Lorsque le payoff se concentre dans une petite région de l'espace des variables motrices~: pour $\theta = 3.8$, l'\omterm{def:qm:estimation:estimator}{erreur type} relative passe de 20~\% à 0.64~\%, soit un facteur d'environ 1\,000 sur la variance.
\textbf{10.} Le rapport de vraisemblance $e^{-\theta Z + \theta^2/2}$ devient très variable~: le moment d'ordre deux $e^{\theta^2}\bar\Phi(z^* + \theta)$ croît de nouveau au-delà de l'optimum.
\textbf{11.} Environ 34 avec les accroissements dans l'ordre temporel et environ 2\,500 avec le \omterm{def:qm:brownian-motion:bridge}{pont brownien}.
\textbf{12.} Les points d'une \omterm{def:qm:monte-carlo:qmc}{suite à discrépance faible} sont les plus uniformes sur leurs premières coordonnées~; le \omterm{def:qm:brownian-motion:bridge}{pont brownien} place là $W_T$ et la forme grossière de la trajectoire, les directions qui portent l'essentiel de la variance du payoff.
\textbf{13.} Environ 37\,000 (\omterm{def:qm:estimation:estimator}{erreur type} 190 fois plus petite)~; pentes ajustées comprises entre $-0.7$ et $-0.8$, contre $-0.5$ pour Monte-Carlo.
\textbf{14.} À partir de la dispersion de brouillages indépendants~: $R$ répliques donnent $R$ \omterm{def:qm:estimation:estimator}{estimateurs} indépendants et sans biais, dont l'écart type empirique divisé par $\sqrt R$ est l'\omterm{def:qm:estimation:estimator}{erreur type}.
\textbf{15.} Ordre fort~: l'erreur trajectorielle, pentes 0.49 (Euler) et 0.89 (Milstein)~; ordre faible~: l'erreur en loi, pente 1.00 pour la moyenne.
\textbf{16.} Coût de $3.2 \times 10^7$ pas de temps contre $3.0 \times 10^9$, soit un facteur 94~: la plupart des tirages se font sur les niveaux grossiers, peu coûteux, et les corrections fines ont une faible variance parce que le \omterm{def:qm:monte-carlo:order}{schéma
de Milstein}, d'ordre fort 1, fait diminuer $V_l$ d'un facteur quatre par niveau.
\textbf{17.} Non~: elle repose sur un payoff régulier, une \omterm{def:qm:monte-carlo:vr}{variable de contrôle} en forme fermée et une faible dimension effective. Les barrières, les digitales et l'exercice anticipé rompent la régularité, et la plupart des transactions n'ont pas de \omterm{def:qm:monte-carlo:vr}{variable de contrôle}
aussi efficace~; chaque famille de transactions doit être mesurée.
\textbf{18.} Compteurs de graine, de transaction et de trajectoire, générateur et version, graines de brouillage, construction des trajectoires et taille de l'ensemble de points.
\textbf{19.} Résultat retenu~: \emph{seize fois moins de machines}~: sur la transaction représentative, la seule \omterm{def:qm:monte-carlo:vr}{variable de contrôle} réduit la variance d'un facteur 853 et, avec les points de Sobol brouillés et le \omterm{def:qm:brownian-motion:bridge}{pont brownien}, d'environ
37\,000, contre les 16 requis~; l'objectif est atteint au nombre actuel de trajectoires, le surcoût de génération restant à mesurer dans le moteur de production.
\textbf{20.} Dans la constante placée devant $n^{-1/2}$~: ce que vaut chaque trajectoire et ce qu'elle coûte.
\end{solution}

\begin{solution}{iq:qm:monte-carlo:1}
L'\omterm{def:qm:estimation:estimator}{erreur type} vaut $\sigma/\sqrt n$ quelle que soit la dimension~: la dimension n'intervient qu'à travers $\sigma$ et le coût d'une trajectoire, ce qui explique que Monte-Carlo l'emporte en grande dimension et perde face aux grilles en une ou deux dimensions.
\iqlookfor{$n^{-1/2}$, indépendant de la dimension, la constante.}
\end{solution}

\begin{solution}{iq:qm:monte-carlo:2}
\omterm{def:qm:monte-carlo:vr}{Variables antithétiques} (payoffs monotones, peu coûteuses, gains faibles)~; \omterm{def:qm:monte-carlo:vr}{variables de contrôle} (une quantité corrélée d'espérance connue~: l'asiatique géométrique, le sous-jacent, un prix de Black--Scholes)~; \omterm{def:qm:monte-carlo:is}{échantillonnage préférentiel}
(événements rares~: déplacer la dérive vers la région qui compte)~; \omterm{def:qm:monte-carlo:vr}{échantillonnage stratifié} ou \omterm{def:qm:monte-carlo:qmc}{quasi-Monte-Carlo} (payoffs réguliers dans quelques directions importantes).
\iqlookfor{Quand chacune fonctionne, et pas seulement son nom.}
\end{solution}

\begin{solution}{iq:qm:monte-carlo:3}
Utiliser un \omterm{def:qm:monte-carlo:prng}{générateur à compteur} (Philox, Threefry) paramétré par la graine, le compteur étant construit à partir des indices de transaction, de trajectoire et de bloc, de sorte que les nombres soient une fonction de leurs coordonnées, et non de l'ordre
de génération~; réduire les sommes partielles dans un ordre fixe.
\iqlookfor{Compteurs plutôt que flux séquentiels, et ordre de réduction fixe.}
\end{solution}

\begin{solution}{iq:qm:monte-carlo:4}
Son uniformité n'est bonne que sur les premières coordonnées, et dans l'ordre temporel la variance de la trajectoire est répartie uniformément sur 250 d'entre elles. Il faut y remédier par un \omterm{def:qm:brownian-motion:bridge}{pont brownien} ou une construction en composantes principales qui
place en premier les directions importantes, et brouiller les points pour obtenir des estimations de l'erreur.
\iqlookfor{Dimension effective et construction des trajectoires.}
\end{solution}

\begin{solution}{iq:qm:monte-carlo:5}
\omterm{def:qm:monte-carlo:is}{Échantillonnage préférentiel}~: déformer la distribution (décaler la moyenne des variables motrices, ou déformer exponentiellement la perte) de sorte que la perte au-delà du seuil devienne courante, et pondérer par le rapport de vraisemblance~; choisir la
déformation au voisinage du point le plus probable de la région rare~; vérifier la dispersion des poids. Le splitting ou le Monte-Carlo conditionnel sont des alternatives.
\iqlookfor{Un \omterm{def:qm:probability-at-speed:change}{changement de probabilité} par le rapport de vraisemblance, et un choix judicieux de la déformation.}
\end{solution}

\begin{solution}{iq:qm:monte-carlo:6}
Écrire l'espérance au niveau le plus fin comme une somme télescopique de corrections entre pas successifs, simuler les trajectoires fine et grossière de chaque correction à partir du même \omterm{def:qm:brownian-motion:bm}{mouvement brownien}, et répartir
les tirages proportionnellement à $\sqrt{V_l/C_l}$. La variance $V_l$ de la correction est gouvernée par l'erreur forte, si bien qu'un ordre fort plus élevé ne demande que très peu de tirages aux niveaux fins~; $V_l$ décroissant
plus vite que le coût ne croît, le coût total est $O(\epsilon^{-2})$, comme pour un \omterm{def:qm:estimation:estimator}{estimateur sans biais}.
\iqlookfor{Somme télescopique, couplage et variance des corrections à partir de l'ordre fort.}
\end{solution}
